% Folder 119, batch 08 — archivists' pages 141–160.
% Pass: Opus 5.5 (claude-opus-5-5), 2026-10-03 — first pass, unchecked against the pages by a human.
%
% Items in this batch (folder 119 special rules, issue #45):
% - pp. 141–142: « Rapport d'activité », undated; his typescript, title in his
%   hand. Already transcribed by the Cercle/CSG (U18d): \page + note only.
% - p. 143: wrapper « Notes de math. oubliées par Jean Malgoire » (probably not
%   his hand); wrapper words given in a note, no \page.
% - pp. 144–145: « Topos quotient d'un espace par une relation d'équivalence
%   locale », dated « Mars 1986 », his hand: transcribed.
% - p. 146: blank divider, no \page.
% - pp. 147–152: « Addendum Esquisse d'un Programme » (his pp. 1–6), undated,
%   his hand: transcribed.
% - p. 153: blank divider, no \page.
% - pp. 154–160: « Digressions combinatoires (pour Réflexions t. 4) », undated,
%   photocopy of his manuscript: transcribed. Disagreement with the survey map:
%   p. 157 is not part of the plan proper — it is a typed list (pairs
%   « Avenir – Passé », « Espace – temps ») with his handwritten renumbering,
%   then his handwritten drafts of footnotes numbered 98–107 for a text of his
%   (referring to « note n° 162 » and « p. 779 »), with sketches of tetrahedra.
%   It is all his own text and hand, so it is transcribed too.
%
\documentclass[11pt,a4paper]{article}
\input{../preamble/grothendieck}

\folder{119}
\batch{8}
\pages{141}{160}
\foldertitle{Esquisse d’un programme [dossier constitué en vue d'une candidature au CNRS (1984)] : copies de tapuscrits et de tapuscrits annotés (1972-1974, 1978-1979, 1984, 1990-1991, s.d.), tirés à part (1971), notes et copies de note manuscrites (1986, s.d.), lettres (1968, 1972, 1991).}
\dating{1968-1991}
\watermark{Édition de démonstration}

\begin{document}

\section{Rapport d'activité}

\page{141}
\note{pp. 141–142 : « Rapport d'activité », tapuscrit non daté de Grothendieck, titre de sa main ; déjà transcrit en ligne par le Cercle/CSG (cote U18d), il n'est pas repris ici. En deux lignes : son enseignement depuis 1978 (C4, DEA, options de DEUG) conçu comme initiation à la recherche, où la difficulté principale des étudiants est la verbalisation ; puis ses directions de thèses (topologie des surfaces) et ses sujets de réflexion — géométrie algébrique anabélienne et tour de Teichmüller, topologie modérée et dévissage des structures stratifiées, fondements de l'algèbre homotopique.}

\page{142}
\note{suite du « Rapport d'activité » ; voir la note de la p. 141.}

\note{p. 143 : chemise portant, d'une main qui n'est probablement pas la sienne, « Notes de math. oubliées par Jean Malgoire » (et « (15 p) »). Elle couvre les pp. 144–160.}

\section{Topos quotient d'un espace par une relation d'équivalence locale}

\page{144}
\marginal{Mettre où ?}

\marginal{Mars 1986 — « Feuilletages »}

\emph{Topos quotient} d'un espace $E$ [ou topos ?] par une relation d'équivalence locale \add{$Q$} (ex : feuilletages de variétés). On trouve comme quotients les espaces de la forme $(X, G)$, $X$ espace, $G$ gr. discret opérant sur $X$ (non libres), $G$ étant le groupe fondamental d'une « feuille » …) [Cas où toutes les feuilles $F_i$ sont homéomorphes, et \ill{} : $E \to F$ induisant des iso $F_i \to F$ …]

\marginal{\uncertain{Notion} \ill{} de l'« holonomie »}

\page{145}
\textbf{Th} Soit $C$ la catégorie des \add{$(X, R)$, $X$} esp. top. \struck{\ill{}} et $R$ rel. d'équivalence locale sur $X$ [localement] ouverte à fibres loc. \struck{\ill{}} connexes, \struck{\ill{}} (flèches : appl. continues compat. avec rel. d'équiv. locales …), et $C'$ la catégorie des \uncertain{flèches} \struck{\ill{}}, systèmes $(X, \mathcal{Y}, f)$, où $X$ espace top., $\mathcal{Y}$ \uncertain{mult. top.}, $f : \mathrm{Top}(X) \to \mathcal{Y}$ morphisme qui est \add{(loct)} ouvert, à fibres loc. connexes et \ill{} (flèches de $C'$ : … \uncertain{i.e.} une 2-cat. dont les $\mathrm{Hom}(\xi, \eta)$ sont des cat. \emph{discrètes}). On a
\[
  C \overset{\approx}{\underset{\approx}{\rightleftarrows}} C' .
\]
Variantes pour variétés \struck{\ill{}} $C^{(r)}$ ($r \in [0, +\infty] \cup \{\omega\} \cup \{\omega_{\mathbb{C}}\}$) feuilletées.

\emph{Question} Quelles variétés obtient-on comme \uncertain{quotients} des feuilletages ???

\marginal{\ill{} \uncertain{cette lacune} \ill{} \uncertain{feuilletages} \ill{}}

\note{la note marginale de gauche, écrite en long, est presque entièrement illisible.}

\section{Addendum Esquisse d'un Programme}

\page{147}
\note{pp. 147–152 : six feuillets numérotés 1 à 6 de sa main.}

Addendum Esquisse d'un Programme.

\emph{Voisinage tubulaire} d'une partie \struck{quelconque} \add{modérée} d'un espace modéré (compact). PS \uncertain{évident} : un point.

I) Sous-espace modéré « bordant » \struck{\ill{}} $B$ d'un espace modéré $X$. Collier autour de $B$
\[
  \bigl(\, \text{\struck{$B \times I \hookrightarrow$}}\; B \subset C \subset X,
  \quad C \text{ vois. de } B,\quad \exists\, C \xrightarrow[\sim]{\varphi} B \times I,
\]
\[
  \text{et } \dot C = \varphi^{-1}(B \times \{1\}) \,\bigr)
\]
\note{il écrit $\varphi(B \times \{1\})$, l'exposant $-1$ n'étant pas visible ; au-dessus de « bordant », « bon “bord” » biffé.}

\struck{$B \simeq B$} avec

\begin{tikzcd}
  & B \arrow[dl] \arrow[dr, "{b \mapsto (b,0)}"] & \\
  C \arrow[rr, "\sim"] & & B \times I
\end{tikzcd}

Collier \struck{\ill{}} \add{strict}. $I = [0, 1]$.

\struck{($(B, 1)$ est bordant dans} ($\dot C$ est bordant dans $X \smallsetminus (C \smallsetminus \dot C)$). Collier rigidifié.

\textbf{Conj} 1) Soit $B$ bordant dans $X$.
\begin{enumerate}
  \item[a)] $\exists$ Collier $(C, \dot C)$, donc aussi collier strict.
  \item[b)] « L'espace » des colliers stricts est contractile.
  \item[c)] L'espace des \struck{colliers \add{rigidifiés} stricts et fibrés sur} collier \struck{\ill{}} est contractile, ou encore : l'espace des autom. d'un collier, induisant l'identité de $B$ (et « \uncertain{libres} » sur $\dot C$) est contractile
\end{enumerate}
i.e. les choix d'un collier \add{strict}, ou d'un collier \add{strict} rigidifié, i.e. d'une rigidification d'un collier, sont anodins.

\marginal{\struck{\ill{}}}

\page{148}
II) Soit $a \in X$. « Cône » autour de $a$ : $a \in C \subset X$, $\dot C \subsetneq C$ tels que
\[
  (C, \dot C) \xrightarrow[\sim]{\rho} \mathrm{Cone}(\text{\struck{$\dot Y$}}\, Y) \qquad (\text{avec } \rho(a) = o).
\]
\note{sous $\dot C$ : « $Y$ \uncertain{espace modéré} ».}

Cône strict : $\dot C$ est bordant dans $X \smallsetminus (C \smallsetminus \dot C)$. Cône rigidifié.

\textbf{Conj}
\begin{enumerate}
  \item[a)] Existe cône \add{strict}
  \item[b)] L'espace des cônes \add{stricts} est contractile
  \item[c)] L'espace des rigidifications des cônes est contractile i.e. l'espace des autom. de $\mathrm{Cone}(Y)$, qui \uncertain{induisent} l'identité sur $Y$ et fixent $o$, est contractile (Alexandroff ?)
\end{enumerate}
Donc le choix d'un cône \struck{strict} autour de $a$, ou d'un cône strict rigidifié, ou d'une \struck{cône} rigidification sur un cône donné, est anodin.

III Catégorie \add{PB} des paires $(X, B)$, $B$ bordant dans $X$.

Catégorie \add{PBCR} des \struck{paires} triples $(X, B, C, \varphi)$, des paires bordantes avec collier \add{strict} rigidifié.

Catégorie PBC $(X, B, C, \dot C)$ (paire bord. avec coll.).

Catégorie \add{EP} des paires $(X, a)$, $X$ pointé

\page{149}
Catégorie \add{EPCR} des triples $(X, a, (C, \varphi))$, $(C, \varphi)$ cône strict \add{rigidifié}.

Catégorie EPC des triples $(X, a, (C, \dot C))$ : espace modéré pointé, \struck{rigidifié} cône strict.

\uncertain{On a} une cascade d'équivalences de cat. \struck{\ill{}} $\infty$-isotopiques

\note{diagramme reproduit avec ses étiquettes ; les deux 2-flèches $\alpha \Rightarrow$ et $\beta \Leftarrow$ entre la verticale « id$_{PB}$ » et les composés latéraux ne sont pas reproduites.}

\begin{tikzcd}[column sep=6em, row sep=large, nodes={font=\scriptsize}]
  & \mathrm{PB} \arrow[dl, "{\text{recoller } B \times I}"'] \arrow[dr, "{\text{recoller Cône}(B)}"] \arrow[dddl, bend right=20, "{\mathrm{id}_{PB}}" description] \arrow[dddr, bend left=20] & \\
  \mathrm{PBCR} \arrow[rr, dashed] \arrow[d, "\text{oubli rig.}"'] & & \mathrm{EPCR} \arrow[d, "\text{oubli rigidification cône}"] \\
  \mathrm{PBC} \arrow[rr, dashed, "{\text{contraction de } B \text{ en un point}}"] \arrow[d, "\text{oubli collier}"'] & & \mathrm{EPC} \arrow[d, "\text{oubli cône}"] \\
  \mathrm{PB} \arrow[rr] & & \mathrm{EP}
\end{tikzcd}

\marginal{à côté de « recoller $B \times I$ » : « $B \times \{0\}$ nouvelle partie bordante » (et « $B \times \{1\}$ »).}

\marginal{Système des 11 foncteurs qui sont des équivalences $\infty$-isotopiques}

\marginal{NB Pour définir $\alpha$ et $\beta$, compliquer les \ill{} \ill{} (i.e. $C$ \ill{}) — PBR}

\struck{\ill{}} \emph{Variante} \struck{Paires} Multipaires bordantes ($B = \coprod_{i \in I} B_i$), espaces multipointés ($I \hookrightarrow X$) ($I$ ens. d'indices fini variable).

\drawing{le même système de catégories redessiné en perspective : PB au sommet, PBCR, EPCR, PBC, EPC, et en bas PB $\to$ EP ; flèches « recoller collier », « recoller cône », « identité », « contraction $B$ » ; un mot biffé \ill{} à gauche.}

\emph{\uncertain{Solution}} : \ill{} \ill{} \ill{} \ill{} $\infty$-isotop., \struck{\ill{}} \ill{} \uncertain{pouvoir} \ill{} travailler avec l'une d'elles, via catégories équivalentes.

\page{150}
IV \emph{Voisinage tubulaire d'une partie} modérée $Y$ d'un espace modéré $X$.

C'est par définition un cône strict \struck{\ill{}} \add{\uncertain{tubulaire}} autour de $o$ dans $X/Y$.

Construction : Soit $f : X \to \mathbb{R}^+$ fonction modérée telle que $f^{-1}(\{0\}) = Y$. On prend $V = f^{-1}([0, \varepsilon])$, $\varepsilon$ petit [et admissible], et $\dot V = f^{-1}(\{\varepsilon\})$.

Rigidification du voisinage tubulaire ($\Leftrightarrow$ du cône). Rigidification $Y$-admissible : j'insère dans le diagramme (\ill{})
\[
  V - Y = V^* \simeq \dot V \times [0, 1[ ,
\]
\[
  X \simeq \bigl( (X \smallsetminus (\dot V \smallsetminus V)) \amalg (\dot V \times I) \bigr) \textstyle\coprod_{\dot V \times \{1\}} Y ,
\]
\note{flèches dessinées de $V - Y$ vers $X$ et de $X$ vers $Y$ ; le recollement porte « $\dot V \times \{1\}$ » en indice ; il écrit bien $X \smallsetminus (\dot V \smallsetminus V)$.}
pour $\varphi : \dot V \times \{1\} \to Y$ appl. modérée (déterminée de façon unique par la rigidification).
\begin{enumerate}
  \item[a)] Rigidification $Y$-admissible existe
  \item[b)] La classe d'homotopie de $\dot V \to Y$ bien déterminée
\end{enumerate}
Mais on voudrait plus précis.

\marginal{Choix de vois. tub. et rigidif. de vois. tub. anodins.}

\marginal{\uncertain{Notons} \ill{} $\dot V$ \ill{}}

\page{151}
1°) Cas où $X$ équisingulier autour de $Y$. Je présume qu'on doit pouvoir prendre $\varphi$ fibrant (« rigidification fibrante »), et qu'on a une équivalence de catégories $\infty$-isotopiques

Syst. $(X', B,$ \struck{$Y$} $Y, \varphi)$ : $(X', B)$ paire bordante, $\varphi : B \to Y$ morphisme fibrant,
\[
  (X', B, Y, \varphi) \longmapsto \bigl( (X' \textstyle\coprod_B (Y, \varphi)),\, Y \bigr)
\]
\[
  \Big\downarrow \approx
\]
\[
  (X, Y) \text{ paires équisingulières.}
\]

\textbf{NB} Les $X \smallsetminus Y$, $Y$ \add{réguliers} \struck{\ill{} alors} doivent correspondre au cas où $X'$ variété avec bord $B$, et $Y$ régulier.

\marginal{\ill{} $Y$ \uncertain{présumé} \ill{} $X - Y$ \uncertain{régulier} \ill{} $Y$ régulier \ill{} ($Y \subset \overline{X - Y}$)}

2°) Cas général (pas nécess. équising.). Il faudrait arriver à imposer des conditions \ill{} sur $\varphi : \dot V \to Y$, pour assurer l'unicité essentielle. Peut être qq chose dans des genres suivants : Pour tout $y \in Y$, si $c(y) = \mathrm{codim}_y(Y, X)$ alors $\dim \varphi^{-1}(y) = c(y) - 1$ (cas des \uncertain{lisses}) plus précis, loct \struck{\ill{}} $\dot V$ est \ill{} sur $Y$

\page{152}
On peut rêver trouver ainsi une description de la catégorie $\infty$-isotopique des paires \add{\uncertain{emplies}} $(X, Y)$, en termes d'une catégorie $(X', B, \varphi : B \to Y)$, où $\varphi$ \struck{serait} soumis à certaines conditions (\uncertain{reflétant} \add{p. ex.} l'exigence dimensionnelle précédente).

\note{trait de séparation.}

Filtrations équidimensionnelles et régulières d'un espace modéré.

Il y en a, je présume, une qui est canonique dans le présent \ill{} de \uncertain{choses}. [\struck{Toute} \add{Plus précis.} \struck{\ill{}} filtration \ill{} \ill{}, il y a une filtration \struck{\ill{}} canonique, régulière équidimensionnelle compatible \ill{} \uncertain{éventuellement}] Les invariants de dévissage de cette filtration sont des invariants de la structure modérée envisagée.

\section{Digressions combinatoires (pour Réflexions t. 4)}

\page{154}
\note{pp. 154–160 : photocopie de son manuscrit. Notation fixée pour ces pages : $\mathrm{Ens}_n$ (ensembles à $n$ éléments), $\mathfrak{S}_n$ pour son $\mathfrak{S}$ gothique (groupes symétriques), $\omega(S)$ (orientations), $S^*$ (dual).}

Digressions combinatoires (pour Réflexions t. 4)

\textbf{Chap 1} Géométrie des ens. de cardinal $\leqslant 4$

\quad NB Faire polygones des \uncertain{sommets}, et dualité \uncertain{rigoureusement} …

\textbf{Chap 2} L'icosaèdre gauche : ses trois \struck{\ill{}} perspectives

\quad \struck{Orientations \ill{}} $\alpha(F) \in \omega(S)$ défini par les trois perspectives

\textbf{Chap 3} Bi-icosaèdres : les trois perspectives. \struck{\ill{}} $\mathrm{Ic}(S) \simeq S^* \times \omega(S) = \mathrm{Bic}(S)$

\textbf{Chap} \struck{\ill{}} Dualité des hexagrammes

\marginal{ou bihexagrammes}

\begin{enumerate}
  \item L'autodualité $\mathrm{Ens}_6 \simeq \mathrm{Ens}_6$ \add{décrite} via $\mathrm{Aut}(S) \simeq \mathrm{Aut}(S^*)$. Application \struck{aux} groupes d'autom. des bi-icosaèdres \struck{\ill{}} \add{et} icos.
  \item Structure bi-icosaédrale par coloriage des arêtes
  \item L'autodualité (trois perspectives)
\end{enumerate}
\begin{enumerate}
    \item[a)] $\forall\, s \in S$, $s^* \in S^*$, $S \smallsetminus \{s\} \simeq S^* \smallsetminus \{s^*\}$
    \item[b)] $R \in$ \struck{\ill{}} $\mathrm{Ar}(S) \times \mathrm{Ar}(S^*)$
    \item[c)] $\mathcal{B}_{3,3}(S) \simeq \mathcal{B}_{3,3}(S^*)$ avec \uncertain{orientations}
\end{enumerate}

4. Le dictionnaire de la dualité

A) \struck{$\mathrm{Ens}_5$} $\mathrm{Ens}_5 \simeq \mathrm{Ens}(1,5) \approx$ Biicos.

\quad $\mathrm{Ens}_5^+ \simeq \mathrm{Ens}^+(1,5) \approx$ Icos

B) $\mathrm{Ens}_2 \times \mathrm{Ens}_4 \approx \mathrm{Ens}_{2,2,2} \approx \mathrm{Cub}_3$

C) $\mathrm{Ens}_{3,3} \approx \mathrm{Ens}_{3,3}$ (ditriades) dualité. Notion bi-ditriades $\equiv$ $\mathrm{Bipunt}_5$

\note{à droite, un premier jet entièrement biffé : « A') Bipunt $\approx$ Biicos. + sommet ; Punt $\approx$ Icos. + sommet ; B') ».}

A') Bi-icos. + sommets $\approx$ \struck{\ill{}}

\quad Icos + — $\approx$ \struck{\ill{}} $\mathrm{Punt}_5$

B') Biicos + ar. $\approx$ \struck{\ill{} $\times$ Car} ; « longueur » \struck{$\mathbb{Z}/2$} $\to$ Car

\note{à droite, un long passage raturé : « autodualité de $\mathrm{Ens}_2 \times \mathrm{Car}$ ».}

C') Biicos + biface $\simeq \mathrm{Ens}_2 \times \mathrm{Ens}_3$ \struck{explication} $\sim$ (hexagrammes)

\quad $\alpha \times f \longleftarrow (\alpha, f)$

\marginal{remonter à Chap 3}

TSVP

\page{155}
\note{la feuille est photocopiée tête-bêche par rapport au numéro des archivistes.}

\struck{B} $B_1$) \struck{\ill{}} $\mathrm{Car}(S) \simeq \mathrm{Car}(S^*)$

\drawing{deux petits carrés, marqués $\varepsilon$, $s(Q)$ et $\varepsilon^*$, $s(Q^*)$, reliés par $\longleftrightarrow$.}
\[
  \varepsilon \wedge \varepsilon^* \wedge \omega(Q) = 1, \qquad \omega(Q) \simeq \omega(Q^*), \qquad
  Q^* \simeq \varepsilon \wedge Q,\quad Q \simeq \varepsilon^* \wedge Q^*
\]
\marginal{explicitement}

$C_1$) Relation avec la donnée d'un \struck{espace} plan affine \struck{$V$} $V$ sur $\mathbb{F}_3$

et donc \struck{card $P(V) = 4$}, $\mathrm{Card}(E = P(V)) = 4$, \struck{Aut}
\[
  \mathfrak{S}_{P(V)} \simeq \mathrm{Aut}(V)(\mathbb{F}_3) \simeq \mathrm{GP}(1, \mathbb{F}_3)
\]
(\uncertain{ou} antipodisme) et une partition $(2,2)$ de $E$, i.e. structure carrée sur $E$.

$\mathrm{card}(E_i) = 3$, $i \in E$ ; $V \xrightarrow{p_i} E_i$, $V \subset \prod E_i$ ; $E_i \times E_j \xleftarrow{\sim} V$ pour $i, j \in E$, $i \neq j$.

\drawing{grille de points.}

Groupe d'automorphismes trilitaire d'ordre $3^3 \times 2^4 = 3 \cdot 144$.

Opère sur $\mathfrak{S}_{\mathcal{E}}$, $\mathcal{E} = \coprod_{i \in E} E_i$, et laisse invariant \struck{trois} un sous-gr. \ill{} des 3 sous-groupes $\mathfrak{S}_\alpha$ ($\alpha \in$ \struck{\ill{}} $\mathrm{Car}(E) = \mathrm{Antip.}(E)$). Le stabilisateur d'un $\mathfrak{S}_\alpha$ est le groupe des automorphismes d'une bi-ditriade. Opère fidèlement sur $\mathfrak{S}_\alpha$, sous-groupe d'ordre 10 de $\mathrm{Aut}(\mathfrak{S}_\alpha)$ (d'ordre $2 \cdot (3!)^2$ \uncertain{)}.

\note{l'« ordre 10 » est lu tel quel ; la lecture « $2 \cdot (3!)^2$ » est incertaine.}

\marginal{les $\mathfrak{S}_\alpha$ sont \ill{} 3-sous-groupes \ill{} Sylow \ill{} \ill{} \ill{} $\mathfrak{S}_\alpha(V_0)$ \ill{} (\ill{})}

\note{la note marginale en biais en bas à gauche, en partie biffée, est presque illisible.}

\page{156}
\note{page de calculs serrée, disposée en tableau ; les étiquettes entourées (1), (2), (3), (1'), (2'), (3'), (1''), (2''), (3'') renvoient aux trois correspondances du dictionnaire de la p. 154.}

(3') $1, 2, 3$ : \drawing{deux cases de deux et trois points} $\mathrm{Ens}_2 \times \mathrm{Ens}_3$ \struck{\ill{}}

(1) $(1, 5) \longleftrightarrow$ (biicosaèdre)

(2) $(2, 4) \leftrightarrow \mathfrak{S}_4 \times \mathfrak{S}_2$ ; $(2, 2, 2)$ i.e. cube (pas \struck{\ill{}} orienté), ext. de $\mathfrak{S}_3$ par $\mathfrak{S}_2^3$

(3) $(3, 3) \circlearrowleft$ ext. de $\mathfrak{S}_2$ par $\mathfrak{S}_3 \times \mathfrak{S}_3$

\quad $\longrightarrow$ avec une 2 structures biicos. i.e. 4 str. icos. formant deux paires d'opposées

\marginal{voir carrés \uncertain{magiques} $(2, 3)$}

$\mathbb{F}_3^2 \simeq \mathrm{Syl}_3$, \quad $D_4 \times \mathbb{F}_2 \simeq \mathrm{Syl}_2$

(2') \drawing{carré à diagonale marquée, avec $\varepsilon$, $a_Q$, $Z_Q$} $\longleftrightarrow$ \drawing{carré à deux diagonales, avec $\varepsilon^*$, $a_Q$, $Z_Q$}
\[
  \omega(Q) \simeq \omega(Q^*) = \omega, \qquad \text{\struck{\ill{}}}\; \boxed{\varepsilon^* \wedge \varepsilon \wedge \omega = 1}
\]
\[
  \boxed{Q^* = \varepsilon \wedge Q^\vee,\quad Q \simeq \varepsilon \wedge Q^*}
\]
\[
  \varepsilon \wedge \mathrm{codiag}\, Q \simeq \varepsilon^* \wedge \mathrm{codiag}\, Q^*
  \qquad (\text{au-dessus : } \omega(S_6),\ \omega(S_6^*))
\]
i.e. $\varepsilon \wedge \varepsilon^* \wedge \mathrm{codiag}\, Q \wedge \mathrm{codiag}\, Q^* \simeq 1$ (avec $\varepsilon \wedge \varepsilon^* = \omega$), d'où $\mathrm{codiag}\, Q^* \simeq \mathrm{diag}\, Q$.

(1'') $1$ + bipentagone $\longleftrightarrow$ $1$ + bipentagone \quad | Biicos. + sommet

\quad le même !

\quad $1$ + pentagone $\longleftrightarrow$ $1$ + pentagone \quad | Icos. + sommet

(2'') $1 \mid 1 \mid$ \drawing{un carré raturé, puis un carré à deux diagonales} $\longleftarrow$ comparer avec (2') \quad | bi-icos. + arête $= \boxed{\times}\ Q^*$, $\varepsilon \simeq \omega(Q^*)$ (*)

\quad | icos. + arête \drawing{petit carré à diagonale, en partie raturé}

(3'') \drawing{deux ovales de 2 et 3 points} \quad comparer avec (3') \quad \struck{$u\ a'\ a\ v\ b\ u'$}

\quad | bi-icos. + biface \drawing{deux triangles hachurés} $\simeq f \times \varepsilon$

\quad | icos. + face, avec pr. dans $f \times \{0, 1\}$

\drawing{icosaèdre en perspective, sommets marqués $u, u', v, v', a, b$, faces hachurées.}
\[
  \{a, b\} \wedge \{u, u'\} \to \{v, v'\}
\]
\[
  \text{\struck{$\{a, b\} \wedge$}} \qquad a \wedge u = v, \qquad a \wedge v = u'
\]

Programme icosaèdre
\[
  \{a, b\} \wedge \{u, u'\} \wedge \{v, v'\} \simeq \text{\struck{\ill{}}}\; \mathrm{diag}\, Q,
  \qquad \{u, u'\} \wedge \{v, v'\} = \mathrm{codiag}(Q),
  \qquad \{a, b\} \simeq \omega(Q)
\]

(*) choix d'une structure icosaédrale dans le biicos-dodé $\simeq \mathrm{diag}(Q)$. Mais ceci
\[
  \simeq \text{\struck{\ill{}}}\; \omega(S \amalg \varepsilon) \simeq \omega(S) \wedge \varepsilon
  \simeq \mathrm{codiag}(Q) \wedge \varepsilon \simeq \mathrm{diag}\, Q \quad \text{OK !}
\]

\page{157}
\note{feuille d'une autre nature intercalée : une liste dactylographiée de couples, renumérotée à la main, puis des projets manuscrits de notes de bas de page numérotées (98 à 107) pour un texte de lui, avec des croquis de tétraèdres ; numéro des archivistes tête-bêche.}

\struck{IV$''_5$} \add{V$''_4$} \struck{Avenir, passé} \add{Devenir}

Avenir — Passé

destinée — histoire

pérennité — ancienneté

innovation — tradition (III)

\struck{le nouveau — l'ancien}

\struck{\ill{}}

\add{élan — enracinement} (II)

\struck{IV$''_6$} \add{V$''_5$} \add{Espace-temps}

Espace — temps

étendue (ou distance) — durée

\struck{\add{ubiquité — éternité}}

\marginal{$= ??$}

\marginal{complètement t. d. m. (+ harmonie) ; mis un peu liste ; fenêtre ouverte ; dessins …}

98 (*) texte barré

\struck{99} (**) Plusieurs existences successives

100. (*) texte barré

(**) Elle qui n'a pourtant \ill{}. Comparer note n° 162, b. d. p. (*) p. 779 « conviction et connaissance »

(***) Réf. cours 1977. Frustration.

103 (*) Autre perspective.

107 (*) Mais pas encore \ill{} la multiplicité \uncertain{\ill{}}.

(**) Copie

(***) \uncertain{chemin} \ill{} \ill{} jusqu'ici par communication

\drawing{plusieurs croquis de tétraèdres et d'octaèdres, en partie hachurés, sommets marqués $a, b, c, c', u$ ; un carré à diagonales ; le nombre 100 en marge.}

\[
  \begin{array}{lcl}
    \text{\struck{$a$}}\,bc \ (\text{entouré}) & & ac \\
    bcu & & bc \\
    ac\,u \ (\text{entouré}) & & ab \\
    \text{délicat}\ abc' & & c'a \\
    \text{\struck{$c'ab$}} \text{ ou } c'ab' & & c'b' \\
    \text{\struck{$c'b$}}u \text{ ou } c'bu' & &
  \end{array}
\]

\page{158}
\emph{Trialité des bi-ditriades}

On peut voir la structure de deux façons, soit par un ens. : 9 éléments $V$ \add{(plan affine sur $\mathbb{F}_3$)}, soit par l'ens. : 12 éléments $E$ (ens. des droites de $V$). Relations d'incidence entre $E$ et $V$.

\marginal{NB Énoncé dans un plan projectif sur $\mathbb{F}_3$ privé d'un point.}

A) La structure sur $V$ revient à la donnée d'une partie $E$ de $\mathcal{B}_3(V)$, telle que
\[
  \begin{cases}
    \mathrm{card}\, E = 12 \\
    \forall\, a, b \in E, \text{ on a } \mathrm{card}\, a \cap b \leqslant 1
  \end{cases}
\]
\struck{\ill{} fait}

\struck{B\ill{}} B) Structure sur $E$ : 1) partition de $E$ en quatre classes de 3 éléments (\uncertain{réunies} les \uncertain{quatre} directions de droites), $E = \coprod_{i \in I} E_i$

2) \struck{Ens.} $V$ de parties : 4 éléments de $E$ (une des droites passant par un point \add{de $V$}, une par projectif) :
\begin{enumerate}
  \item[a)] Tout \struck{\ill{}} $D \in V$ rencontre chaque $E_i$ en un pt \add{(au moins)} (\struck{\ill{}} \add{au plus}, \struck{\ill{}} \add{revient au même}) \struck{(donc $D \to$)} i.e. $D \to I$ est bijectif)
  \item[b)] Si $D, D' \in V$, $D \neq D'$, alors $\mathrm{Card}\, D \cap D' = 1$ \quad \textbf{NB} b) résulte de a) et c) : \ill{}
  \item[c)] Si $a, b \in E$ n'appartiennent pas au même $E_i$, $\exists$ unique $D \in V$ qui y passe. (unicité provient d'ailleurs de b))
\end{enumerate}
\textbf{NB} c implique déjà que $\mathrm{card}\, V = 9$, et compte tenu de a), que par tout point $a \in E$ passent exact 3 « droites ».

\marginal{$E \cong V \amalg \{\infty\}$ ; \ill{} \ill{} et $E$ \ill{} \ill{} de $C$ \ill{} \ill{} (\ill{}) dans $\mathrm{Ens}$ \ill{} : chaque \uncertain{droite} \ill{} $(E, \dot V)$ \ill{} \ill{} \ill{}}

\note{la note marginale de gauche, en biais, est en grande partie illisible.}

\page{159}
Géométrie des ens. finis (de card $\leqslant$ 3) (Plan)

I Ens : deux éléments
\begin{enumerate}
  \item Produit contracté — $\mathfrak{S}_2 \simeq \pm 1 \simeq \mathbb{Z}/2\mathbb{Z} \simeq \mathbb{F}_3^*$ (deux interprétations)
  \item \struck{O\ill{}} Orientation d'un ens. fini ($\omega(E)$). Associativité ; lien avec orientation des vectoriels réels, ou des vectoriels orthogonaux \add{(car $\neq 2$)}, ou les vectoriels sur $\mathbb{F}_3$.
  \item Hypercubes et antipodismes. Orientation \uncertain{induite}
  \item Carrés (sans lien avec carrés et cubes ordinaires)
\end{enumerate}

II Ens : trois éléments
\begin{enumerate}
  \item 2-vectoriels sur $\mathbb{F}_2$ (ou \add{droite} \struck{\ill{}} projective sur $\mathbb{F}_2$) : $\mathfrak{S}_3 \simeq \mathrm{Gl}(2, \mathbb{F}_2) \simeq \mathrm{Sl}(2, \mathbb{F}_2) \simeq \mathrm{GP}(1, \mathbb{F}_2)$
  \item Droite affine sur $\mathbb{F}_3$. Application aux ditriades et biditriades. $\mathfrak{S}_3 \simeq \mathrm{Aff}(1, \mathbb{F}_3)$ ($\to \mathbb{F}_3^* \simeq \pm 1$) ;
  \[
    1 \to \mathfrak{S}_3 \times \mathfrak{S}_3 \to \mathrm{Aut} \text{ d'une bi-ditriade} \to \text{\struck{\ill{}}}\; 1
  \]
  \item \struck{Quadritriades (ou tri-biditriades)}
\end{enumerate}

III Ens : quatre éléments
\begin{enumerate}
  \item \struck{Espace} \add{Plan} affine sur $\mathbb{F}_2$ : $\mathfrak{S}_4 \simeq \mathrm{Aff}(2, \mathbb{F}_2)$
  \item Lien avec cubes : $\mathrm{Ens}_4 \times \mathrm{Ens}_2 \approx \mathrm{Cub}_3$ ; $\mathfrak{S}_4 \simeq \mathrm{Aut}^+(\mathrm{Cub}_3) \simeq$ \struck{\ill{}} $\mathfrak{S}_{3,4}$ \uncertain{autom.} d'une \uncertain{partition} de deux 3 sur $\mathbb{F}_3$ qui \uncertain{inversent} une sym. des trois droites formant les bords \ill{}
  \item Lien avec \uncertain{V} \ill{} automorphismes
  \item Droite projective sur $\mathbb{F}_3$ : $\mathfrak{S}_4 \simeq \mathrm{GP}(1, \mathbb{F}_3)$
  \item Les quadritriades (ou \struck{tri-bi}triades).
  \[
    1 \to (\mathfrak{S}_3 \times \mathfrak{S}_3)' \to \mathrm{Aut} \text{ d'une quadritriade} \to \mathfrak{S}_4 \to 1,
  \]
  \[
    (\mathfrak{S}_3 \times \mathfrak{S}_3)' = \{(u, v) \mid \varepsilon(u) = \varepsilon(v)\},
    \qquad \mathrm{Aut} \text{ d'une quadritriade} \simeq \mathrm{Aff}(2, \mathbb{F}_3)
  \]
\end{enumerate}

IV Ens : 5 éléments
\begin{enumerate}
  \item \add{Droites proj. sur corps : 4 droites} ($\mathfrak{S}_5^+ \simeq \mathrm{GP}(1, \mathbb{F}_4)$) \quad $\mathfrak{S}_5 \simeq$ Aut. droite projective sur un corps : 4 droites non précisé

  (Si $E$ droite proj. sur $k \simeq \mathbb{F}_4$, alors $\omega(E) \simeq (k^* \smallsetminus \{1\}) =$ ens. des deux gén. de $k$ sur $\mathbb{F}_2$.)

  Explication \uncertain{renvoyant} aux \uncertain{anharmoniques}
  \item $\mathfrak{S}_5 \simeq \mathrm{GP}(1, \mathbb{F}_5)$ sera développé avec bi-icosaèdre
\end{enumerate}

\note{en III.2 la fin de la ligne (« autom. d'une … bords ») est une addition marginale serrée dont la lecture est très incertaine ; en IV.1 la ligne « Droites proj. sur corps : 4 droites » est écrite au-dessus et entourée avec la formule.}

\page{160}
V Théorie combinatoire des ens. : 6 élmts et des bi-icosaèdres [Peut être scindé en deux : A) \struck{I}Icosaèdre et bi-icosaèdre \quad B) géométrie des ens. à 6 éléments]

Voir plus détaillé ailleurs …

\struck{\ill{}}

\emph{Appendice} Polygones combinatoires.

Pour bien faire, il faudrait un exposé sur les relations géométriques entre les configurations envisagées ici …

\end{document}
